In this section, we introduce in detail our static
program analysis algorithm, {\THESYSTEM} for estimating the \emph{adaptivity} of a program,
which corresponds to
%  as the adaptivity the 
third part at the bottom of Figure~\ref{fig:adaptfun-structure}. 
% It effectively combines the quantitative and reachability analysis techniques for analyzing
% adaptivity by considering its two sub-properties.
% The effective combination complements the weakness of two lines of work in comparison with any single one.
% Also, it computes more accurate adaptivity comparing to the previous adaptivity analysis with higher efficiency.
It is presented through the following sections.

\paragraph{{Chapter Outline}}
%
In order to give a sound and precise approximation of the formal adaptivity from Section~\ref{sec:adapt-exe} efficiently, 
the new static program analysis framework is developed through four steps analysis in two phases
similar to the semantics-based adaptivity formalization procedure.
\begin{itemize}
   \item The first phase estimate the \emph{semantics-based} dependency graph mainly in two steps.
   \begin{enumerate}
   \item In Section~\ref{sec:alg_edgegen}, we first analyze the \emph{reachability} sub-property
   %  --
   % the data \emph{dependency relation} 
   through a combined data and control flow analysis technique, and estimate the
   \emph{variable may-dependency} relation.
   \item In Section~\ref{sec:alg_weightgen}, we analyze the \emph{quantitative} sub-property
   %  -- the \emph{dependency depth} 
   through
   the reachability bound analysis techniques, and estimates the weight of each labeled variable on the \emph{semantics-based dependency graph}.
   \item Then in Section~\ref{sec:alg_graphgen}, we construct the \emph{estimated dependency graph} encoding two estimated results above, together with the statically collected vertices and query annotation.
   \end{enumerate}
   \item The second phase in Section~\ref{sec:alg_adaptcompute} computes program's adaptivity.
   % the program adaptivity estimation, 
   In this phase, we construct a program-based dependence graph for approximating the execution-based graph in Section~\ref{sec:dynamic-adapt}.
   Then, based on this graph, we design an algorithm based on graph theory for
   computing the adaptivity upper bound soundly 
   and accurately.
   \end{itemize}